\SourceLeaf{98}%
imaginärer, durch eine unendliche Ebene veranschaulichen, in der jeder durch die Abszisse $=a$ und die Ordinate $=b$ bestimmte Punkt gleichsam die Größe $a+bi$ darstellt. Der stetige Übergang von einem Wert von $x$ zu einem anderen Wert $a+bi$ geschieht demnach entlang einer Linie und ist somit auf unendlich viele Arten möglich. Ich behaupte nun, dass das Integral $\int \varphi x\mathord{.} dx$ auf zwei verschiedenen Wegen immer denselben Wert erhält, wenn in dem von den beiden Weglinien eingeschlossenen Flächenbereich nirgends $\varphi x = \infty$ gilt. Das ist ein sehr schöner Lehrsatz\HistoricalNote{\textsuperscript{*})}{*)}{Genau genommen wird hierbei noch vorausgesetzt, dass $\varphi x$ selbst eine einwertige Funktion von $x$ ist oder dass man zumindest im gesamten Flächenbereich ein einziges System ihrer Werte ohne Unterbrechung der Stetigkeit zugrunde legt.}, dessen nicht besonders schwierigen Beweis ich bei passender Gelegenheit geben werde. Er hängt mit anderen schönen Wahrheiten über Reihenentwicklungen zusammen. Der Übergang zu jedem Punkt ist stets möglich, ohne je eine Stelle zu berühren, an der $\varphi x=\infty$ gilt. Ich verlange daher, dass man solche Punkte umgeht, an denen der ursprüngliche Grundbegriff von $\int \varphi x\mathord{.} dx$ offensichtlich seine Klarheit verliert und leicht zu Widersprüchen führt. Daraus wird zugleich klar, wie eine durch $\int \varphi x\mathord{.} dx$ erzeugte Funktion für denselben Wert von $x$ mehrere Werte haben kann. Beim Übergang dorthin kann man nämlich einen Punkt, an dem $\varphi x=\infty$ gilt, entweder gar nicht, einmal oder mehrmals umlaufen. Definiert man zum\,Beispiel $\log x$ durch $\int \frac{1}{x}\,dx$, beginnend bei $x=1$, so gelangt man zu $\log x$, ohne den Punkt $x=0$ einzuschließen oder indem man ihn einmal oder mehrmals umläuft. Jedes Mal kommt dann die Konstante $+2\pi i$ oder $-2\pi i$ hinzu. Damit sind die verschiedenen Logarithmen einer jeden Zahl ganz klar. Kann $\varphi x$ für keinen endlichen Wert von $x$ unendlich werden, so ist das Integral stets eine einwertige Funktion. Dies ist zum\,Beispiel für $\varphi x = \dfrac{e^x-1}{x}$ der Fall, sodass $\int \dfrac{e^x-1}{x}\,dx$ gewiss eine einwertige Funktion von $x$ ist. Ihr Wert wird durch die \textls{immer} konvergierende Reihe dargestellt, die also stets einen und nur \textls{einen} Sinn hat,
\[
x+\tfrac{1}{4}xx+\tfrac{1}{18}x^3+\tfrac{1}{96}x^4+\text{etc.}
\]
\noindent Ich wünschte, Herr \textsc{Soldner} hätte, da er eine neue Funktion%
